% Source printed pp. 54-59.
\section{Castelnuovo--Mumford 正则性}\label{src:2.3}
\begin{epigraph}
代数几何中一些最深刻的结果, 涉及给出层的高阶上同调群为 $0$ 的判别准则的问题...我们将在这里(借助 Nakai, Matsusaka 和 Kleiman 发展并使用的技术)只证明一个较弱的消灭定理, 但它能统一地适用于一大类层.
\epigraphauthor{David Mumford~\cite[第 $99$ 页]{Mum66a}}
\end{epigraph}
Serre 消灭定理---有时称为 Cartan--Serre--Grothendieck 定理---断言: 若 $F$ 为射影簇 $(X,\mathcal O_X(1))$ 上的凝聚层, 则当 $d\gg0$ 时, 
\begin{enumerate}[(1)]
\item $F(d)$ 无基点; 
\item 对 $i>0$, $H^i(X,F(d))=0$; 且
\item 对所有 $k\geq0$, 乘法映射
\[
H^0(X,F(d))\otimes H^0(X,\mathcal O(k))
\longrightarrow H^0(X,F(d+k))
\]
满射.
\end{enumerate}
Castelnuovo--Mumford 正则性定量衡量了使扭转 $F(d)$ 满足这三个上同调性质所需的 $d$ 的大小.
\subsection{定义与基本性质}\label{src:2.3.1-subsection}
\begin{definition}\label{src:2.3.1}
设 $F$ 为域 $k$ 上射影空间 $\mathbb P^n$ 的凝聚层.
对整数 $m$, 称 $F$ 为 \newterm{$m$-正则的}{m-regular}, 如果
\[
H^i(\mathbb P^n,F(m-i))=0
\]
对所有 $i\geq1$ 成立.
$F$ 的\newterm{正则性}{regularity}是使 $F(m)$ 为 $m$-正则的最小整数 $m$.
\end{definition}
由正则性的定义可知, 若 $F$ 为 $m$-正则, 则 $F(d)$ 为 $(m-d)$-正则.
我们将在命题~\bookref{2.3.5}~中证明: 若 $F$ 为 $m$-正则, 则对所有 $d\geq m$, 它也是 $d$-正则的.
定义中要求第 $(m-i)$ 次扭转的第 $i$ 上同调消灭, 乍看可能有些神秘, 但这个定义对于按维数归纳的论证很方便, 以下结果就说明了这一点.
\begin{lemma}\label{src:2.3.2}
设 $F$ 为域 $k$ 上 $\mathbb P^n$ 的 $m$-正则凝聚层.
若超平面 $H\subseteq\mathbb P^n$ 避开 $F$ 的伴随点, 则 $F|_H$ 也是 $m$-正则的.
\end{lemma}
\begin{proof}
假设表明, 在仿射开子概形 $U\subseteq\mathbb P^n$ 上, $H$ 的定义方程为模 $\Gamma(U,F)$ 的非零因子.
因此 $H:F(-1)\to F$ 为单射, 且对每个整数 $i>0$, 有短正合列
\[
0\longrightarrow F(m-i-1)\longrightarrow F(m-i)
\longrightarrow F|_H(m-i)\longrightarrow0.
\]
它诱导上同调长正合列
\[
\cdots\longrightarrow H^i(\mathbb P^n,F(m-i))
\longrightarrow H^i(H,F|_H(m-i))
\longrightarrow H^{i+1}(\mathbb P^n,F(m-i-1))\longrightarrow\cdots.
\]
若 $F$ 为 $m$-正则, 则
$H^i(\mathbb P^n,F(m-i))=H^{i+1}(\mathbb P^n,F(m-i-1))=0$.
所以对所有 $i>0$, $H^i(H,F|_H(m-i))=0$, 故 $F|_H$ 也是 $m$-正则的.
\end{proof}
\begin{exercise}[简单, 适合练习]\label{src:2.3.3}
\begin{enumerate}[(a)]
\item 证明 $\mathbb P^n$ 上的 $\mathcal O(d)$ 为 $(-d)$-正则.
\item 证明次数为 $d$ 的超曲面 $H\subseteq\mathbb P^n$ 的结构层为 $(d-1)$-正则.
\item 证明亏格为 $g$ 的光滑曲线 $C\subseteq\mathbb P^n$ 的结构层为 $(2g-1)$-正则.
\end{enumerate}
\end{exercise}
\begin{exercise}[简单]\label{src:2.3.4}
设 $0\to K\to F\to Q\to0$ 为 $\mathbb P^n$ 上凝聚层的短正合列.
若 $K$ 为 $(m+1)$-正则而 $F$ 为 $m$-正则, 证明 $Q$ 为 $m$-正则.
\end{exercise}
正则性的另一个优点是如下 Castelnuovo 的结果.
\begin{proposition}[正则性的性质]\label{src:2.3.5}
设 $F$ 为域 $k$ 上 $\mathbb P^n$ 的 $m$-正则凝聚层.
\begin{enumerate}[(1)]
\item 对 $d\geq m$, $F$ 为 $d$-正则.
\item 若 $d\geq m$ 且 $k\geq0$, 则乘法映射
\[
H^0(\mathbb P^n,F(d))\otimes H^0(\mathbb P^n,\mathcal O(k))
\longrightarrow H^0(\mathbb P^n,F(d+k))
\]
满射.
\item 对 $d\geq m$, $F(d)$ 整体生成, 且对 $i\geq1$, $H^i(\mathbb P^n,F(d))=0$.
\end{enumerate}
\end{proposition}
\begin{proof}
若 $k\to k'$ 为域扩张, 则平坦基变换(\bookref{A.2.12})表明
\[
H^i(\mathbb P_k^n,F)\otimes_k k'
=H^i(\mathbb P_{k'}^n,F\otimes_k k').
\]
由于 $k\to k'$ 忠实平坦, 断言~(1)--(3)~可在基变换后检验.
故可假设 $k$ 代数闭, 特别地, 为无限域.
对~(1)~和~(2), 我们按 $n$ 归纳; 初始情形 $n=0$ 显然成立.
若 $n>0$, 由于 $k$ 无限, 可以选择避开 $F$ 的伴随点的超平面 $H\subseteq\mathbb P^n$.
限制 $F|_H$ 在 $H\cong\mathbb P^{n-1}$ 上为 $m$-正则(引理~\bookref{2.3.2}), 故归纳假设表明~(1)~和~(2)~对 $F|_H$ 成立.

为证明~(1), 还按整数 $d$ 归纳. 初始情形 $d=m$ 由假设成立.
当 $d>m$ 时, 短正合列 $0\to F(d-i-1)\to F(d-i)\to F|_H(d-i)\to0$ 诱导上同调长正合列
\[
\cdots\longrightarrow H^i(\mathbb P^n,F(d-i-1))
\longrightarrow H^i(\mathbb P^n,F(d-i))
\longrightarrow H^i(H,F|_H(d-i))\longrightarrow\cdots.
\]
当 $i>0$ 时, 第一项由关于 $d$ 的归纳假设消灭, 因为 $F$ 为 $(d-1)$-正则.
第三项由关于 $n$ 的归纳假设消灭: 由引理~\bookref{2.3.2}, $F|_H$ 为 $m$-正则, 因而由关于 $n$ 的归纳假设为 $d$-正则, 故 $H^i(H,F|_H(d-i))=0$.
于是第二项消灭, 得证~(1).

为证明~(2), 除 $n$ 外, 我们还按 $k$ 归纳.
将乘法映射记为
\[
\mu_{d,k}:H^0(\mathbb P^n,F(d))\otimes H^0(\mathbb P^n,\mathcal O(k))
\longrightarrow H^0(\mathbb P^n,F(d+k)).
\]
虽然初始情形 $k=0$ 显然成立, 但归纳论证需要直接证明 $k=1$ 的情形.
为此, 考虑交换图
\begin{equation}\label{src:2.3.6}
\begin{tikzcd}[column sep=large,row sep=large,font=\small]
& H^0(\mathbb P^n,F(d))\otimes H^0(\mathbb P^n,\mathcal O(1))
\arrow[r,"\nu_d\otimes\mathrm{res}"]\arrow[d,"\mu_{d,1}"]&
H^0(H,F|_H(d))\otimes H^0(H,\mathcal O_H(1))\arrow[d]\\
H^0(\mathbb P^n,F(d))\arrow[r,"\alpha"]
\arrow[ur,"\mathrm{id}\otimes H"]&
H^0(\mathbb P^n,F(d+1))\arrow[r,"\nu_{d+1}"]&
H^0(H,F|_H(d+1)).
\end{tikzcd}
\end{equation}
映射 $\alpha$ 由乘以 $H\in H^0(\mathbb P^n,\mathcal O(1))$ 给出, 且有包含 $\operatorname{im}(\alpha)\subseteq\operatorname{im}(\mu_{d,1})$.
由于由~(1)~有 $H^1(\mathbb P^n,F(d))=0$, 限制映射 $\nu_d:H^0(\mathbb P^n,F(d))\to H^0(H,F|_H(d))$ 满射.
同样, 由于 $H^1(\mathbb P^n,\mathcal O)=0$, $\mathrm{res}:H^0(\mathbb P^n,\mathcal O(1))\to H^0(H,\mathcal O_H(1))$ 满射, 故上方水平箭头满射.
将归纳假设应用于 $H=\mathbb P^{n-1}$, 可知右方竖直箭头满射.
因此, 复合 $\nu_{d+1}\circ\mu_{d,1}$ 满射, 从而 $\operatorname{im}(\mu_{d,1})$ 满射到 $H^0(H,F|_H(d+1))$.
由下行的正合性, 
\[
\begin{aligned}
H^0(\mathbb P^n,F(d+1))
&=\operatorname{im}(\mu_{d,1})+\ker(\nu_{d+1})\\
&=\operatorname{im}(\mu_{d,1})+\operatorname{im}(\alpha)
=\operatorname{im}(\mu_{d,1}),
\end{aligned}
\]
故 $\mu_{d,1}$ 满射.

若 $k>1$, 考虑交换方形
\[
\begin{tikzcd}[column sep=normal,row sep=large,font=\footnotesize]
H^0(\mathbb P^n,F(d))\otimes H^0(\mathbb P^n,\mathcal O(k-1))
\otimes H^0(\mathbb P^n,\mathcal O(1))
\arrow[r]\arrow[d,"\mu_{d,k-1}\otimes\mathrm{id}"']&
H^0(\mathbb P^n,F(d))\otimes H^0(\mathbb P^n,\mathcal O(k))
\arrow[d,"\mu_{d,k}"]\\
H^0(\mathbb P^n,F(d+k-1))\otimes H^0(\mathbb P^n,\mathcal O(1))
\arrow[r,"\mu_{d+k-1,1}"]&
H^0(\mathbb P^n,F(d+k)).
\end{tikzcd}
\]
分别对 $k-1$ 和 $k=1$ 应用归纳假设, 左方竖直映射和下方水平箭头均满射.
故 $\mu_{d,k}$ 满射.

为证明~(3), 我们知道当 $k\gg0$ 时, $F(d+k)$ 整体生成, 即
$\gamma_{F(d+k)}:H^0(\mathbb P^n,F(d+k))\otimes\mathcal O_{\mathbb P^n}\to F(d+k)$
满射.
考虑交换方形
\[
\begin{tikzcd}[column sep=normal,row sep=large,font=\footnotesize]
H^0(\mathbb P^n,F(d))\otimes H^0(\mathbb P^n,\mathcal O(k))
\otimes\mathcal O_{\mathbb P^n}
\arrow[r,"\mu_{d,k}\otimes\mathrm{id}"]
\arrow[d,"\gamma_{F(d)}\otimes\mathrm{id}"']&
H^0(\mathbb P^n,F(d+k))\otimes\mathcal O_{\mathbb P^n}
\arrow[d,"\gamma_{F(d+k)}"]\\
F(d)\otimes\bigl(H^0(\mathbb P^n,\mathcal O(k))\otimes\mathcal O_{\mathbb P^n}\bigr)
\arrow[r,"\mathrm{id}\otimes\gamma_{\mathcal O(k)}"]&
F(d)\otimes\mathcal O(k).
\end{tikzcd}
\]
由~(2), 上方水平箭头满射, 故从左上到右下的复合满射.
由于下方水平映射由与 $F(d)$ 作张量积定义, 映射 $\gamma_{F(d)}$ 必须满射.
最后, 为说明 $F(d)$ 的高阶上同调消灭, 注意对每个 $i>0$, 由~(1), 层 $F$ 为 $(d+i)$-正则, 故 $H^i(\mathbb P^n,F(d))=0$.
\end{proof}
断言~(1)~的一个直接推论是: 若 $F$ 为 $m$-正则, 则对所有 $d\geq m$, 限制映射
\[
\nu_d::H^0(\mathbb P^n,F(d))\longrightarrow H^0(H,F|_H(d))
\]
满射. 事实上, (1)~表明 $F$ 也为 $d$-正则, 满射性来自 $H^1(\mathbb P^n,F(d-1))$ 的消灭.
以下引理将在定理~\bookref{2.3.8}~的证明中使用, 它表明在较弱的假设下仍可保证 $\nu_d$ 满射.
\begin{lemma}\label{src:2.3.7}
设 $F$ 为 $\mathbb P^n$ 上的凝聚层, $H$ 为避开 $F$ 的伴随点的超平面.
若 $F|_H$ 为 $m$-正则, 且 $\nu_d$ 对某个 $d\geq m$ 满射, 则 $\nu_p$ 对所有 $p\geq d$ 满射.
\end{lemma}
\begin{proof}
观察图~\eqref{src:2.3.6}, 可知上方箭头 $\nu_d\otimes\mathrm{res}$ 满射(因为 $\nu_d$ 和 $\mathrm{res}$ 均满射), 右方竖直乘法态射也满射(对 $m$-正则层 $F|_H$ 应用命题~\bookref{2.3.5}(2)).
这意味着 $\nu_{d+1}$ 满射, 断言由归纳法成立.
\end{proof}
\subsection{正则性的界}\label{src:2.3.2-subsection}
我们证明平凡向量丛的子层的正则性有界, 这一结果由 Mumford 在~\cite[第 $101$ 页]{Mum66a}~中首次建立.
它是代数几何中几乎每个模空间有界性的基础.\footnote{即使在~\cite[Tag 0DPA]{SP}~中使用 Artin 代数性公理(\bookref{C.7.4})构造 Hilbert 概形与 Quot 概形时, 建立拟紧性所用的也是这一结果~\cite[Tag 08AG]{SP}.}
\begin{theorem}[正则性的有界性]\label{src:2.3.8}
对每一对非负整数 $r,n$ 及每个多项式 $P\in\mathbb Q[z]$, 存在整数 $m_0$, 满足如下性质: 对每个域 $k$, 具有 Hilbert 多项式 $P$ 的每个子层 $F\subset\mathcal O_{\mathbb P_k^n}^{\oplus r}$ 都为 $m_0$-正则.
\end{theorem}
\begin{proof}
由于 $m$-正则性不受域扩张影响, 可以假设 $k$ 无限.
我们按 $n$ 归纳. 初始情形 $n=0$ 成立, 因为 $\mathbb P^0$ 上每个层 $F$ 对每个整数 $m$ 都是 $m$-正则的.

设 $n\geq1$, $F\subseteq\mathcal O_{\mathbb P^n}^{\oplus r}$ 具有 Hilbert 多项式 $P$.
可以选择避开 $\mathcal O_{\mathbb P^n}^{\oplus r}/F$ 的所有伴随点的超平面 $H\subseteq\mathbb P^n$.
这保证
$\operatorname{Tor}^{\mathcal O_{\mathbb P^n}}_1
(\mathcal O_H,\mathcal O_{\mathbb P^n}^{\oplus r}/F)=0$, 
且短正合列 $0\to F\to\mathcal O_{\mathbb P^n}^{\oplus r}\to\mathcal O_{\mathbb P^n}^{\oplus r}/F\to0$
限制为短正合列
\begin{equation}\label{src:2.3.9}
0\longrightarrow F|_H\longrightarrow\mathcal O_H^{\oplus r}
\longrightarrow\mathcal O_H^{\oplus r}/F\longrightarrow0.
\end{equation}
由于 $H\cong\mathbb P^{n-1}$, 可以对 $F|_H\subseteq\mathcal O_H^{\oplus r}$ 应用归纳假设.
另一方面, 由于 $F\subseteq\mathcal O_{\mathbb P^n}^{\oplus r}$ 无挠, 有短正合列
\begin{equation}\label{src:2.3.10}
0\longrightarrow F(-1)\xrightarrow{H}F\longrightarrow F|_H\longrightarrow0.
\end{equation}
$F|_H$ 的 Hilbert 多项式为
\[
\chi(F|_H(d))=\chi(F(d))-\chi(F(d-1))=P(d)-P(d-1).
\]
特别地, $F|_H$ 的 Hilbert 多项式仅依赖于 $P$.
对 $F|_H\subseteq\mathcal O_H^{\oplus r}$ 应用归纳假设, 得到整数 $m_1$, 使 $F|_H$ 为 $m_1$-正则.

对 $m\geq m_1$, 由于 $F|_H$ 为 $m$-正则, 故也为 $(m+1)$-正则, 因此对所有 $i\geq1$, 
\[
H^i(H,F|_H(m-i))=H^{i+1}(H,F|_H(m-(i+1)))=0.
\]
将~\eqref{src:2.3.10}~扭转 $m-i$ 后所诱导的上同调长正合列, 表明
\[
H^{i+1}(\mathbb P^n,F(m-(i+1)))
\xrightarrow{\sim}H^{i+1}(\mathbb P^n,F(m-i)).
\]
由于 $H^{i+1}(\mathbb P^n,F(d))$ 在 $d\gg0$ 时消灭, 可以断言: 对 $m\geq m_1$, $F$ 满足上同调次数至少为 $2$ 的 $m$-正则性条件.

为处理 $H^1$, 注意当 $m\geq m_1-1$ 时, 长正合列
\begin{equation}\label{src:2.3.11}
\begin{aligned}
0&\longrightarrow H^0(\mathbb P^n,F(m-1))
\longrightarrow H^0(\mathbb P^n,F(m))
\longrightarrow H^0(H,F|_H(m))\\
&\longrightarrow H^1(\mathbb P^n,F(m-1))
\longrightarrow H^1(\mathbb P^n,F(m))
\longrightarrow H^1(H,F|_H(m))=0
\end{aligned}
\end{equation}
给出不等式 $h^1(\mathbb P^n,F(m_1-2))\geq h^1(\mathbb P^n,F(m_1-1))\geq\cdots$, 它们最终稳定为 $0$.
我们断言, 事实上, 在达到 $0$ 之前, 不等式
$h^1(\mathbb P^n,F(m_1-2))>h^1(\mathbb P^n,F(m_1-1))>\cdots$
都是严格的.
为说明这一点, 注意当 $m\geq m_1-1$ 时, 
$h^1(\mathbb P^n,F(m-1))=h^1(\mathbb P^n,F(m))$
当且仅当 $\nu_m:H^0(\mathbb P^n,F(m))\to H^0(H,F|_H(m))$ 满射.
假设对某个 $m'\geq m$, $\nu_{m'}$ 满射.
由于 $F|_H$ 为 $m_1$-正则, 可用引理~\bookref{2.3.7}~推出: 对所有 $m''\geq m'$, $\nu_{m''}$ 满射.
因此 $h^1(\mathbb P^n,F(m''))$ 对 $m''\geq m'$ 为常数, 故为零.
断言得证. 令 $m_2=m_1+h^1(\mathbb P^n,F(m_1))$, 则
$h^1(\mathbb P^n,F(m_2-1))=0$, 故 $F$ 为 $m_2$-正则.

还需证明 $m_2$ 被一个不依赖于 $F$ 的常数 $m_0$ 控制.
由于 $F\subseteq\mathcal O_{\mathbb P^n}^{\oplus r}$, 对任意 $d\geq0$, 
$h^0(\mathbb P^n,F(d))\leq rh^0(\mathbb P^n,\mathcal O(d))=r\binom{n+d}{n}$.
利用 $i\geq2$ 时 $h^i(\mathbb P^n,F(m_1))$ 的消灭, 得到
\[
\begin{aligned}
h^1(\mathbb P^n,F(m_1))
&=h^0(\mathbb P^n,F(m_1))-\chi(F(m_1))\\
&\leq r\binom{n+m_1}{n}-P(m_1).
\end{aligned}
\]
定义 $m_0:=m_1+1+r\binom{n+m_1}{n}-P(m_1)\geq m_2$, 
就得到所需的常数, 使具有 Hilbert 多项式 $P$ 的每个 $F\subseteq\mathcal O_{\mathbb P^n}^{\oplus r}$ 均为 $m_0$-正则.
\end{proof}
\begin{remark}\label{src:2.3.12}
上述证明实际上建立了一个更强的断言.
为陈述这一结果, 回忆: 每个次数为 $n$ 的数值多项式 $P\in\mathbb Q[z]$(即对整数 $d\gg0$, 有 $P(d)\in\mathbb Z$)均可唯一写成
\[
P(d)=\sum_{i=0}^na_i\binom di,\qquad a_i\in\mathbb Z.
\]
这可由直接的归纳论证得到, 参见~\cite[命题 I.7.3]{Har77}.
对非负整数 $r,n$, 存在多项式
\begin{equation}\label{src:2.3.13}
\Lambda_{r,n}\in\mathbb Z[x_0,\ldots,x_n],
\end{equation}
满足如下性质: 对每个域 $k$, 具有 Hilbert 多项式
$P(d)=\sum_i a_i\binom di$ 的每个子层
$F\subseteq\mathcal O_{\mathbb P_k^n}^{\oplus r}$
均为 $m_0$-正则, 其中 $m_0=\Lambda_{r,n}(a_0,\ldots,a_n)$.
见~\cite[第 $101$ 页]{Mum66a}.
\end{remark}
\begin{exercise}\label{src:2.3.14}
设 $P=d$ 为固定整数 $d$ 对应的常数多项式.
使每个定义长度为 $d$ 的 $0$ 维子概形的理想层 $I\subseteq\mathcal O_{\mathbb P^n}$ 为 $m$-正则的最优 $m$ 是多少?
\end{exercise}
\begin{remark}[最优界]\label{src:2.3.15}
虽然 Mumford 的正则性有界性结果(定理~\bookref{2.3.8})给出了显式界, 且足以用于包括构造 Hilbert 概形在内的许多应用, Gotzmann 建立了更优的界: 
对域 $k$ 上具有 Hilbert 多项式 $P$ 的射影概形 $X\subseteq\mathbb P^N$, 存在唯一的整数
$\lambda_1\geq\lambda_2\geq\cdots\geq\lambda_r\geq1$, 使 $P$ 可写成
\[
P(d)=\binom{d+\lambda_1-1}{\lambda_1-1}
+\binom{d+\lambda_2-2}{\lambda_2-1}
+\cdots+\binom{d+\lambda_r-r}{\lambda_r-1},
\]
且 $X$ 的理想层 $I_X$ 为 $r$-正则.
见~\cite{Got78}, \cite[\S$4.3$]{BH93}, \cite{Gre89}~和~\cite[\S$3$]{Gre98}.
\end{remark}
\begin{exercise}\label{src:2.3.16}
设 $C\subseteq\mathbb P^n$ 为次数 $d$, 亏格 $g$ 的曲线.
证明 Gotzmann 的界表明 $C$ 的理想层 $I_C$ 为 $\bigl(\binom d2+1-g\bigr)$-正则.
能否将它与定理~\bookref{2.3.8}~的证明给出的界比较, 即能否从~\eqref{src:2.3.13}~计算 $\Lambda_{1,n}(1-g,d)$?
\end{exercise}
\begin{remark}\label{src:2.3.17}
\cite{GLP83}~证明了: 次数为 $d$ 的整的非退化曲线 $C\subseteq\mathbb P^N$ 的理想层 $I_C$ 为 $(d-N+2)$-正则.
更一般的猜想断言: 维数 $n$, 次数 $d$ 的光滑非退化射影子簇 $X\subseteq\mathbb P^N$ 的理想层为 $(d-(N-n))+1)$-正则; 见~\cite{GLP83,EG84}.
\end{remark}
\begin{proposition}[族中的正则性]\label{src:2.3.18}
设 $S$ 为 Noether 概形, $\pi:\mathbb P_S^n\to S$ 为相对射影空间, $Q$ 为 $\mathbb P_S^n$ 上在 $S$ 上平坦的凝聚层.
假设存在整数 $m>0$, 使对每个 $s\in S$, 在 $\mathbb P_{\kappa(s)}^n$ 上的限制 $Q_s$ 都为 $m$-正则.
则对 $d\geq m$, 
\begin{enumerate}[(1)]
\item $\pi_*Q(d)$ 为向量丛, 且其构造与沿 $f:T\to S$ 的基变换交换, 即 $f^*\pi_*Q(d)\cong\pi_{T,*}Q_T(d)$; 
\item 对 $i>0$, $R^i\pi_*Q(d)=0$; 且
\item $\pi^*\pi_*Q(d)\to Q(d)$ 满射.
\end{enumerate}
\end{proposition}
\begin{proof}
由于对所有 $s\in S$, 
$R^i\pi_*Q(d)\otimes\kappa(s)\to H^i(\mathbb P_{\kappa(s)}^n,Q_s(d))=0$
满射, 上同调与基变换(\bookref{A.6.8})给出~(1)~和~(2).
对~(3), 令 $C=\operatorname{coker}(\pi^*\pi_*Q(d)\to Q(d))$.
对每个 $s\in S$, 有 Cartesian 图
\[
\begin{tikzcd}
\mathbb P_{\kappa(s)}^n\arrow[r,"j"]\arrow[d,"\pi_s"']
\arrow[dr,phantom,"\square"]&
\mathbb P_S^n\arrow[d,"\pi"]\\
\operatorname{Spec}\kappa(s)\arrow[r,"i"]&S.
\end{tikzcd}
\]
利用~(1), 有
$j^*\pi^*\pi_*Q(d)=\pi_s^*i^*\pi_*Q(d)=\pi_s^*\pi_{s,*}Q_s(d)$.
伴随映射 $\pi^*\pi_*Q(d)\to Q(d)$ 沿 $j$ 的拉回, 对应于伴随映射
$\pi_s^*\pi_{s,*}Q_s(d)\to Q_s(d)$; 
由正则性的性质(\bookref{2.3.5}(3)), 后者满射.
因此, 对所有 $s\in S$, 
$C\otimes\kappa(s)=\operatorname{coker}(\pi_s^*\pi_{s,*}Q_s(d)\to Q_s(d))=0$, 
故 $\pi^*\pi_*Q(d)\to Q(d)$ 满射.
\end{proof}
